\scanpage{162}
\mspage{107}
de l'axiome pertinent d'associativité de la composition
et de celui sur les unités, mais seulement que
les $1_x$ soient des unités bilatères pour la
composition.

En supposant donné \unsure{ce|un} supplément de
structure sur le foncteur $\mathrm{Ch}\deleted{\unsure{(X)|(\ )}}$, \deleted{\illegible{2}} et
revenant à la considération du foncteur
associé $C:\mathcal M_{/S}\longrightarrow\underline{\mathrm{Fib}}\,S$,
on trouve un\deleted{\unsure{e} application} morphisme canonique
\[
C(C(X))\xrightarrow{\gamma_X}C(X)
\]
par
\[
(x,c,c')\longmapsto(x,c'\circ c),
\]
et l'axiome des unités bilatères nous
assure aussitôt que
\[
\unsure{(*)|(a)|(\alpha)}\qquad\qquad \gamma_X\circ\deleted{\unsure{p}}i_{C(X)}=\gamma_X\circ C(i_X)=\mathrm{id}_{C(X)}.
\]
Comme $i_{C(X)}\in W$, cela prouve déjà que
$\gamma_X\in W_{\deleted{S}}$, donc $\gamma_X\in W_S$. Donc $\gamma_X$ devient
un iso dans $(\underline{\mathrm{Fib}}\,S)W_S^{-1}$, et \deleted{\unsure{on aura}}\inserted{la}
relation $\unsure{(*)|(a)|(\alpha)}$ implique bien $i_{C(X)}=C(i_X)$
dans ce localisé, q.e.d.

Je vais à nouveau résumer ce que j'ai
dégagé.
% archive sheet number: 161
